\documentclass{standalone}
\usepackage{circuitikz}
\usetikzlibrary{calc}
\definecolor{circuitnotemark}{HTML}{FE00FE}
\begin{document}
\begin{circuitikz}[american, line width=0.8pt]
\ctikzset{bipoles/length=1.2cm}
\foreach \x in {0,2,4,6,8,10,12,14,16} {\foreach \y in {0,-2,-4} {\fill[gray, opacity=0.35] (\x,\y) circle (1.1pt);}}
\node[gray, font=\scriptsize] at (0,0.84) {1};
\node[gray, font=\scriptsize] at (2,0.84) {2};
\node[gray, font=\scriptsize] at (4,0.84) {3};
\node[gray, font=\scriptsize] at (6,0.84) {4};
\node[gray, font=\scriptsize] at (8,0.84) {5};
\node[gray, font=\scriptsize] at (10,0.84) {6};
\node[gray, font=\scriptsize] at (12,0.84) {7};
\node[gray, font=\scriptsize] at (14,0.84) {8};
\node[gray, font=\scriptsize] at (16,0.84) {9};
\node[gray, font=\scriptsize, anchor=east] at (-0.84,0) {a};
\node[gray, font=\scriptsize, anchor=east] at (-0.84,-2) {b};
\node[gray, font=\scriptsize, anchor=east] at (-0.84,-4) {c};
\coordinate (b2) at (2,-2);
\coordinate (b7) at (12,-2);
\coordinate (a1) at (0,0);
\coordinate (c1) at (0,-4);
\coordinate (b4) at (6,-2);
\coordinate (a6) at (10,0);
\coordinate (c6) at (10,-4);
\coordinate (a9) at (16,0);
\coordinate (c9) at (16,-4);
\node[plain amp] (part-U1) at (b2) {}; % line 2
\node[font=\scriptsize, anchor=north] at (part-U1.south) {$\mathrm{LM358}$}; % line 2
\draw ($(part-U1.+)+(0.44,-0.13)$) -- ++(0.28,0); % line 2
\draw ($(part-U1.+)+(0.58,-0.27)$) -- ++(0,0.28); % line 2
\draw ($(part-U1.-)+(0.44,0.13)$) -- ++(0.28,0); % line 2
\node[transformer] (part-T1) at (b7) {}; % line 3
\node[font=\scriptsize, anchor=north] at (part-T1.south) {$\mathrm{1to1}$}; % line 3
\draw (a1) -| (part-U1.+); % line 5
\draw (c1) -| (part-U1.-); % line 6
\draw (part-U1.out) -| (b4); % line 7
\draw (a6) -| (part-T1.A1); % line 8
\draw (c6) -| (part-T1.A2); % line 9
\draw (a9) -| (part-T1.B1); % line 10
\draw (c9) -| (part-T1.B2); % line 11
\path (0,-6.395) rectangle (3.048,-5.605); % line 13
\node[anchor=west, inner sep=0, circuitnotemark, font=\footnotesize] at (0,-6) {X}; % line 13
\path (10,-6.395) rectangle (13.895,-5.605); % line 14
\node[anchor=west, inner sep=0, circuitnotemark, font=\footnotesize] at (10,-6) {X}; % line 14
\end{circuitikz}
\end{document}
